\documentclass[12pt,a4paper]{amsart}

\usepackage{amsmath,amssymb,amsthm}
\usepackage{mathtools}
\usepackage[colorlinks=true,linkcolor=blue!60!black,citecolor=green!50!black,
            urlcolor=blue!70!black]{hyperref}
\usepackage[shortlabels]{enumitem}
\usepackage{booktabs}

%% ---- Theorem environments ----------------------------------------
\newtheorem{theorem}{Theorem}[section]
\newtheorem{lemma}[theorem]{Lemma}
\newtheorem{proposition}[theorem]{Proposition}
\newtheorem{corollary}[theorem]{Corollary}
\theoremstyle{definition}
\newtheorem{definition}[theorem]{Definition}
\newtheorem{remark}[theorem]{Remark}
\newtheorem{assumption}[theorem]{Assumption}
\newtheorem{conjecture}[theorem]{Conjecture}
\newtheorem{problem}{Problem}[section]

%% ---- Notation (aligned with Paper I) ----------------------------
\newcommand{\PW}{\mathrm{PW}_\omega}
\newcommand{\Kop}{\mathcal{K}_c}
\newcommand{\GVM}{G^{(N)}_{\mathbf{p},c}}
\newcommand{\QM}{Q^{(N)}_{\mathbf{p},c}}
\newcommand{\psin}{\psi_n^{(c)}}
\newcommand{\EN}{E_{N,c}}
\newcommand{\norm}[1]{\|#1\|}
\newcommand{\ip}[2]{\langle #1,\, #2\rangle}
\newcommand{\Rbulk}{\mathcal{R}_{\mathrm{bulk}}}
\newcommand{\Rturning}{\mathcal{R}_{\mathrm{turn}}}
\newcommand{\PN}{P_{N,c}}
\newcommand{\CXRY}{C_{\mathrm{XRY}}}

\title[Conditional Frame Stability via PSWF Spectral Projection]
      {Conditional Frame Stability of Prolate Sampling Operators\\
       via PSWF Spectral Projection and XRY Quadrature:\\
       A Conditional Implication Framework for Quadrature Compatibility}

\thanks{This paper is a companion to \cite{PaperI}, where frame coercivity
and a conditional banded decay estimate were established under the
Discrete Spectral Transfer Property (DSTP) of \cite{PaperI}.
For consistency with \cite{PaperI}, the compatibility condition referred
to as Assumption~3.10 in earlier drafts of this paper is now called the
DSTP throughout.
The present paper does not verify the DSTP outright. Instead, it
organises the problem into two explicit conjectural components: a PSWF
product spectral tail conjecture and an XRY spectral stability conjecture.
Their combination yields a conditional verification mechanism for the
DSTP used in \cite{PaperI}.
The PSWF product tail conjecture (Conjecture~\ref{conj:pswf_product_tail})
is now known to be false as a global uniform statement;
the edge obstruction is documented in \cite{PaperIII}.}

\author{Anonymous Author}
\address{}
\email{}
\date{\today}

\subjclass[2020]{42C05, 42B10, 47B32, 41A60, 65D32}

\keywords{prolate spheroidal wave functions, PSWF-Gauss quadrature,
finite frames, frame stability, Gram operator, spectral projection,
quadrature compatibility, Xiao--Rokhlin--Yarvin, exponential convergence,
edge obstruction, bandwidth doubling}

\begin{document}
\maketitle

\begin{abstract}
In the companion paper \cite{PaperI}, frame coercivity and a conditional
banded decay estimate for the quadrature defect matrix $\EN$ were
established for the prolate sampling operator on $V_{N,c}$, conditional
on the Discrete Spectral Transfer Property (DSTP) of \cite{PaperI}.

For consistency with \cite{PaperI}, we refer to the compatibility
condition (Assumption~3.10 in earlier drafts) as the DSTP throughout
this paper.

The present paper does not verify the DSTP unconditionally.
Instead, for PSWF-Gauss quadrature rules of Xiao--Rokhlin--Yarvin
\cite{Xiao2001}, it organises the compatibility problem into two explicit
conjectural inputs. First, one needs a PSWF product spectral tail statement
for the analytic product family
$f_{mn}=\psi_m^{(c)}\overline{\psi_n^{(c)}}$ in the relevant index regime.
Second, one needs an XRY spectral stability statement on the same family.
Assuming these two conjectural inputs in the bulk and off-diagonal regimes,
one obtains the conditional defect estimate
\[
  |E_{mn}| \le C e^{-\alpha N} \qquad (0\le m,n\le N-1,\; m,n\le\gamma N).
\]
The edge regime $m,n\sim N$ is identified as an obstruction:
no uniform exponential bound of this form can hold globally,
as proved in \cite{PaperIII} via a bandwidth doubling argument.
Consequently, the DSTP of \cite{PaperI} is conditionally verified
in the bulk and off-diagonal regimes under the stated inputs.
\end{abstract}

\tableofcontents

%% ==================================================================
\section{Introduction}
%% ==================================================================

\subsection*{Context}

The companion paper \cite{PaperI} established the following structure
for the quadrature defect matrix $\EN = \GVM - G^{\mathrm{ref}}_{N,c}$:
\begin{enumerate}
\item An exact identity $E_{mn} = R_{mn}^{\mathrm{quad}}$
  (pure quadrature remainder; PSWF orthogonality term vanishes exactly).
\item A conditional banded decay estimate: under the DSTP of
  \cite{PaperI},
  \[
    |E_{mn}| \le C_p(1+|m-n|)^{-p} + Cc^{-1/3}
    \qquad (0 \le m,n \le N-1).
  \]
\item A conditional frame stability theorem with explicit bounds.
\end{enumerate}
The open problem left by \cite{PaperI} was to identify a concrete
quadrature class satisfying the DSTP.

\subsection*{Contribution of this paper}

This paper does not solve the full construction problem unconditionally.
Its contribution is an implication framework: for XRY PSWF-Gauss quadrature,
the DSTP of \cite{PaperI} follows from two separate conjectural
inputs, namely a PSWF product spectral tail conjecture and an XRY spectral
stability conjecture on the relevant product family
$f_{mn}=\psi_m^{(c)}\overline{\psi_n^{(c)}}$, restricted to the bulk
and off-diagonal regimes.
The companion paper \cite{PaperIII} has since established that a
globally uniform exponential bound is false at the edge ($m,n\sim N$)
due to bandwidth doubling; see Conjecture~\ref{conj:pswf_product_tail}(iii)
and Remark~\ref{rem:edge-failure}.

\subsection*{Strategy of this paper}

The technical strategy consists of an implication chain with two explicit
analytical inputs:
\begin{enumerate}[(i)]
\item \emph{PSWF product spectral tail conjecture.} The product integrands
  $f_{mn}=\psi_m^{(c)}\overline{\psi_n^{(c)}}$ are expected to admit
  sufficiently uniform PSWF spectral tail decay in the bulk and
  off-diagonal regimes. The edge regime is an obstruction.
\item \emph{XRY spectral stability conjecture.} The XRY quadrature
  functional $Q_N$ is expected to satisfy a stability bound on the same
  product family, controlled by the PSWF spectral truncation error
  associated with $\PN$.
\end{enumerate}
Within the present framework, these two inputs are used only through their
combination in the deduction of the defect estimate. The technical content
of the paper is therefore the exact implication chain from these two
conjectural inputs to the DSTP of \cite{PaperI}, together
with the perturbative passage from that condition to frame stability.

\subsection*{Main result}

\begin{theorem}[Conditional implication theorem]\label{thm:main_informal}
Assume parts~(i) and~(ii) of Conjecture~\ref{conj:pswf_product_tail}
and Conjecture~\ref{conj:xry_stability},
together with the auxiliary technical hypotheses stated in
Assumptions~\ref{ass:pswf_sup}, \ref{ass:node_sep}, and~\ref{ass:weights},
restricted to the bulk regime $m,n\le\gamma N$.
Then for PSWF-Gauss quadrature rules \cite{Xiao2001}, the quadrature defect
satisfies $|E_{mn}|\le Ce^{-\alpha N}$ for all $0\le m,n\le\gamma N$.
Consequently, the DSTP of \cite{PaperI} holds conditionally
in the bulk and off-diagonal regimes, and the frame stability theorem
of \cite{PaperI} applies to PSWF-Gauss sampling in those regimes.
The edge regime $m,n\sim N$ is excluded; see part~(iii) of
Conjecture~\ref{conj:pswf_product_tail} and \cite{PaperIII}.
\end{theorem}

The formal version is Theorem~\ref{thm:compatibility_verified} together
with Theorem~\ref{thm:frame_stability}.

\subsection*{Positioning}

This paper should be read as an implication framework (in the sense of a
reduction between explicit analytic hypotheses) rather than a Theorem-paper
in the traditional sense. Its purpose is to isolate the two core unresolved
analytical components that stand between the abstract framework of
\cite{PaperI} and a concrete verification of the DSTP for XRY PSWF-Gauss
quadrature.
The PSWF product-tail step and the XRY stability step are both left as
explicit conjectural inputs. In particular, the present work is not a
classical analysis paper deriving the conclusion from weaker previously
established estimates; rather, it makes explicit where the unresolved
analytic burden lies and how the downstream frame-theoretic conclusion would
follow once those inputs are available.
Paper~III \cite{PaperIII} has since resolved the edge question negatively
and has established unconditional first-moment spectral localization
in the off-diagonal regime; these results are incorporated into
Conjecture~\ref{conj:pswf_product_tail} and Remarks~\ref{rem:bottleneck}
and~\ref{rem:edge-failure} below.

\subsection*{Organisation}

Section~\ref{sec:pswf_gauss} sets up PSWF-Gauss quadrature and the
spectral projector. Section~\ref{sec:spectral_decay} states the PSWF-side
conjectural input and explains the scope of the product-tail step.
Section~\ref{sec:xry_accuracy} states the XRY spectral stability conjecture.
Section~\ref{sec:verification} derives the conditional verification of
the DSTP. Section~\ref{sec:stability} states the conditional frame
stability theorem. Section~\ref{sec:open} records open problems.

%% ==================================================================
\section{PSWF-Gauss Quadrature and Spectral Projector}\label{sec:pswf_gauss}
%% ==================================================================

\begin{definition}[PSWF spectral projector]\label{def:projector}
Let $\{\psi_k^{(c)}\}_{k\ge0}$ be an orthonormal PSWF basis in
$L^2([-T,T])$. For $N\ge1$, define
\[
  \PN f := \sum_{k=0}^{N-1}\ip{f}{\psi_k^{(c)}}\,\psi_k^{(c)}.
\]
\end{definition}

\begin{definition}[PSWF-Gauss quadrature operator]\label{def:xry}
Let $(\mathbf{p},\mathbf{w})=((p_1,\dots,p_N),(w_1,\dots,w_N))$ denote
the PSWF-Gauss nodes and weights of Xiao--Rokhlin--Yarvin \cite{Xiao2001}.
The quadrature functional is $Q_N(f):=\sum_{j=1}^N w_j f(p_j)$.
\end{definition}

\begin{remark}[No algebraic closure claim]\label{rem:no_algebra}
We do not assume any closure of the PSWF system under pointwise
multiplication. All arguments below use only analytic regularity and
spectral approximation properties.
\end{remark}

\begin{assumption}[Node separation]\label{ass:node_sep}
The PSWF-Gauss nodes satisfy
$\min_{j\ne k}|p_j-p_k|\ge c_0 T/N$
for some constant $c_0>0$ in the regime $N\sim c$.
\end{assumption}

\begin{assumption}[Weight positivity and boundedness]\label{ass:weights}
The weights satisfy $w_j>0$ and $C_1 T/N\le w_j\le C_2 T/N$
uniformly in $j$ in the regime $N\sim c$.
\end{assumption}

\begin{remark}[Node-weight assumptions and evidence]\label{rem:nodes_weights_evidence}
Assumptions~\ref{ass:node_sep} and~\ref{ass:weights} are included here as
explicit hypotheses rather than derived facts. They are, however,
consistent with the numerical behaviour reported for PSWF-Gauss rules in
\cite{Xiao2001,OsipovRokhlinXiao}, where the nodes appear well separated and
the weights remain positive and of the expected order $T/N$ in the regime
$N\sim c$. We do not claim these properties as proved here.
\end{remark}

%% ==================================================================
\section{PSWF Product Spectral Tail Conjecture}\label{sec:spectral_decay}
%% ==================================================================

\begin{assumption}[Uniform PSWF sup-norm control]\label{ass:pswf_sup}
There exists a constant $C_{\infty}>0$ such that
\[
  \norm{\psi_n^{(c)}}_{L^\infty([-T,T])}\le C_{\infty}
\]
uniformly for $0\le n\le N\sim c$ in the bulk regime; compare the standard
PSWF bounds discussed in \cite{Slepian1978,OsipovRokhlinXiao}.
\end{assumption}

\begin{lemma}[Uniform product norm bound]\label{lem:product_norm}
Assume Assumption~\ref{ass:pswf_sup} and the regime $0\le m,n\le N\sim c$.
Then there exists a constant $C>0$ such that
\[
  \norm{f_{mn}}_{L^2([-T,T])}
  = \norm{\psi_m^{(c)}\overline{\psi_n^{(c)}}}_{L^2([-T,T])}
  \le C
\]
uniformly in $m,n$.
\end{lemma}
\begin{proof}
By Assumption~\ref{ass:pswf_sup},
$\norm{\psi_m^{(c)}}_{L^\infty([-T,T])}\le C_{\infty}$ uniformly in the
bulk regime, while $\norm{\psi_n^{(c)}}_{L^2([-T,T])}=1$ by orthonormality.
Hence
$\norm{f_{mn}}_{L^2}
  \le \norm{\psi_m^{(c)}}_{L^\infty}\norm{\psi_n^{(c)}}_{L^2}
  \le C_{\infty}.$
\end{proof}

\begin{remark}[Role of the product norm bound]\label{rem:generic_bound}
The preceding estimate uses only the generic $L^\infty\cdot L^2$ bound and
carries no PSWF-specific spectral information. Its role is solely to ensure
that the $L^2$ factor in Conjecture~\ref{conj:xry_stability} remains
uniformly bounded; the nontrivial spectral content is entirely contained in
Conjecture~\ref{conj:pswf_product_tail}.
\end{remark}

\begin{remark}[Why the PSWF product-tail issue is difficult]
The family $f_{mn}=\psi_m^{(c)}\overline{\psi_n^{(c)}}$ is not known to enjoy
any finite-dimensional closure property in the PSWF basis. Thus a uniform
spectral-tail statement for $0\le m,n\le N\sim c$ is not a routine corollary
of standard PSWF approximation results, but a separate open analytical
problem --- and one that is now known to have a negative answer at the edge;
see part~(iii) below.
\end{remark}

\begin{conjecture}[PSWF product spectral tail --- revised]\label{conj:pswf_product_tail}
For the family
$\mathcal{F}_N:=\{\psi_m^{(c)}\overline{\psi_n^{(c)}}:0\le m,n\le N\}$
and the regime $N\sim c$:

\begin{enumerate}[(i)]
\item \emph{Bulk regime.}
For $m,n\le \gamma N$ with fixed $\gamma<1$,
there exist constants $C_{\mathrm{b}}(\gamma),\alpha_{\mathrm{b}}(\gamma)>0$
such that
\[
  \|(I-\PN)f_{mn}\|_{L^2([-T,T])}
  \le C_{\mathrm{b}}\,e^{-\alpha_{\mathrm{b}} N}
\]
uniformly in $m,n\le\gamma N$.
This is the regime used in the conditional implication of this paper.
A proof route via exponential convolution decay of PSWF Fourier tails
is identified as Problem~\ref{prob:bulkconj}.

\item \emph{Off-diagonal regime.}
For $|m-n|\ge\delta N$ with fixed $\delta>0$,
there exist $C_p(\delta)>0$ for every $p\ge1$ such that
\[
  \|(I-\PN)f_{mn}\|_{L^2([-T,T])}
  \le C_p(\delta)\,(1+|m-n|)^{-p}
\]
uniformly over $0\le m,n\le N$.
Conditional pointwise spectral localization supporting this
estimate is Assumption~3.1 of \cite{PaperIII}.
An unconditional uniform $L^2$ bound
$\|(I-\PN)f_{mn}\|\le C T^{1/2}$ holds without assumptions
(Theorem~1 of \cite{PaperIII}); the algebraic-decay
form above requires additional structural input.

\item \emph{Edge obstruction (proved).}
No uniform exponential bound of the form in~(i) can hold for indices
approaching the edge $m,n\sim N$.
Precisely: for $n$ in the Landau--Widom transition zone
$|n - 2c/\pi|\le c^{1/2}$, the bandwidth doubling of
$\widehat{f_{nn}}=\hat\psi_n^{(c)}*\hat\psi_n^{(c)}$ implies
\[
  \mathcal{E}_{\mathrm{out}}(f_{nn}) \ge c_0 > 0,
\]
so the Slepian bound cannot yield exponential tail decay
(Proposition~4.1 of \cite{PaperIII}).
This is the fundamental obstruction to a globally uniform statement
and is not a conjecture but a proved fact.
\end{enumerate}

The form actually used in the conditional implication of this paper
is part~(i). Parts~(ii) and~(iii) record the current state
of knowledge and the limits of the global conjecture.
\end{conjecture}

\begin{remark}[Edge failure is structural, not technical]\label{rem:edge-failure}
Part~(iii) is not a limitation of the proof method.
Bandwidth doubling forces $\operatorname{supp}\widehat{f_{mn}}\subset[-2\omega,2\omega]$,
so the out-of-band energy $\mathcal{E}_{\mathrm{out}}(f_{mn})$
measures the Fourier mass of $f_{mn}$ in $(\omega,2\omega)$.
For edge indices this mass is bounded away from zero uniformly in $N$.
Consequently, any argument based on the Slepian energy identity fails
at the edge, regardless of how refined the analysis is.
The edge regime must be handled --- if at all --- by entirely different methods.
\end{remark}

\begin{remark}[Refined bottleneck: from absence of structure to
explicit energy gap]\label{rem:bottleneck}
The companion paper \cite{PaperIII} has established two unconditional
results that replace the previously open structural question with
a precisely identified analytic gap:
\begin{enumerate}[(a)]
\item \emph{Mean spectral localization} (Proposition~3.2 of
  \cite{PaperIII}): for all $0\le m,n\le N$,
  \[
    \mathbb{E}_{mn}[\chi_k] = \mu_{mn} + E_{mn},
  \]
  where $\mu_{mn}=\chi_m(c)+\chi_n(c)$ and $E_{mn}$ is given
  by an explicit energy formula.
\item \emph{Spectral lower bound} (Proposition~3.3 of \cite{PaperIII}):
  \[
    \mathbb{E}_{mn}[\chi_k] \ge \mu_{mn}/2
  \]
  unconditionally for all $0\le m,n\le N$ with $f_{mn}\not\equiv0$.
\end{enumerate}
The obstruction is therefore no longer the absence of intermediate
structure, but rather the lack of sufficiently sharp control on the
weighted energy terms in the commutator identity:
specifically, whether the kinetic and potential contributions average
to approximately $\mu_{mn}/2$ under the product measure $\psi_m^2\psi_n^2\,dx$.
This is an energy equidistribution problem with a semiclassical character;
see Problem~4.2 of \cite{PaperIII}.
\end{remark}

\begin{remark}[Unconditional vs.\ conditional tail bounds]\label{rem:unconditional-gap}
Paper~III \cite{PaperIII} provides the following hierarchy
of unconditional results in the off-diagonal regime ($|m-n|\ge\delta N$):
\begin{center}
\renewcommand{\arraystretch}{1.3}
\begin{tabular}{lll}
\toprule
Bound & Status & Source \\
\midrule
$\|(I-\PN)f_{mn}\|\le CT^{1/2}$ & \textbf{unconditional} & Thm.~1 of \cite{PaperIII} \\
$\mathbb{E}_{mn}[\chi_k]\ge\mu_{mn}/2$ & \textbf{unconditional} & Prop.~3.3 of \cite{PaperIII} \\
$\|(I-\PN)f_{mn}\|\le C_p(1+|m-n|)^{-p}$ & conditional & Thm.~2 of \cite{PaperIII} \\
\bottomrule
\end{tabular}
\end{center}
The unconditional $T^{1/2}$ bound does not exhibit decay in $|m-n|$
and is therefore insufficient for the exponential defect estimate
required by the DSTP of \cite{PaperI} without the additional structural
input of Conjecture~\ref{conj:pswf_product_tail}(ii).
Any quantitative improvement toward algebraic or exponential decay
requires either the semiclassical energy equidistribution result
(Problem~4.2 of \cite{PaperIII}) or a direct spectral localization
assumption.
\end{remark}

\begin{remark}[Frequency-space heuristic for the product-tail problem]
PSWFs with bandlimit parameter $c$ may be viewed heuristically as functions
whose Fourier content is essentially concentrated on a frequency window of
size comparable to $c$. Under this heuristic, the pointwise product
$\psi_m^{(c)}\overline{\psi_n^{(c)}}$ corresponds in frequency space to a
convolution, so one expects an effective bandwidth of order $2c$ rather
than $c$. Since the projector $\PN$ retains only the first $N\sim c$ PSWF
modes, the main issue is not whether spectral mass extends beyond that
range, but how rapidly the corresponding PSWF coefficients decay past the
cutoff. Part~(iii) of Conjecture~\ref{conj:pswf_product_tail} shows that
this decay fails at the edge.
\end{remark}

\begin{remark}[Analytical bottleneck --- updated]\label{rem:bottleneck-implication}
The overall implication chain shows that the main unresolved analytical
difficulty is concentrated in the PSWF product spectral tail in the
bulk regime $m,n\le\gamma N$ (Conjecture~\ref{conj:pswf_product_tail}(i))
and the off-diagonal algebraic-decay regime
(Conjecture~\ref{conj:pswf_product_tail}(ii)).
The edge regime is excluded by the proved obstruction in part~(iii).
Once the bulk tail estimate is available, the remaining steps
reduce to quadrature stability and perturbative frame analysis.
\end{remark}

%% ==================================================================
\section{XRY Spectral Stability Conjecture}\label{sec:xry_accuracy}
%% ==================================================================

\begin{remark}[Why the XRY stability conjecture is structurally plausible]
The XRY quadrature rule is constructed independently through the
PSWF-Gauss node-weight scheme of Xiao--Rokhlin--Yarvin \cite{Xiao2001}, and
not by building the projector $\PN$ into the quadrature functional itself.
The role of $\PN$ in Conjecture~\ref{conj:xry_stability} is therefore
an a posteriori analytical comparison device rather than part of the
construction of $Q_N$. Moreover, the XRY formulas are exact on the
underlying PSWF trial space, so $Q_N$ integrates $\PN f$ exactly whenever
$\PN f\in V_{N,c}$. This suggests that $Q_N$ behaves as an approximate
spectral projector onto $V_{N,c}$ at the level of integration functionals,
with the defect governed by the contribution of $(I-\PN)f$. With that
distinction in place, it is reasonable to expect the quadrature error to be
governed primarily by the PSWF spectral tail of the test function, up to a
residual term that is exponentially small in the regime $N\sim c$. This is
also consistent with the interpolation formulas underlying the XRY
construction. This does not prove Conjecture~\ref{conj:xry_stability}, but
it explains why a bound of the form
\[
  |Q_N(f)-\textstyle\int_{-T}^T f|\lesssim \|(I-\PN)f\| + e^{-\alpha N}\norm{f}_{L^2}
\]
represents a natural stability statement rather than an ad hoc
reformulation of the original compatibility problem.
\end{remark}

\begin{conjecture}[XRY stability conjecture on the PSWF product family]
\label{conj:xry_stability}
Let
$\mathcal{F}_N:=\{\psi_m^{(c)}\overline{\psi_n^{(c)}}:0\le m,n\le N\}$.
There exist constants $\CXRY,C,\alpha>0$, depending on the XRY
construction parameters and on the regime $N\sim c$, such that for every
$f\in\mathcal{F}_N$,
\[
  \left|Q_N(f)-\int_{-T}^T f(x)\,dx\right|
  \le \CXRY\,\|(I-\PN)f\|_{L^2([-T,T])}
      + C\,e^{-\alpha N}\,\norm{f}_{L^2([-T,T])}.
\]
This conjecture formalises the XRY-side stability input needed in the
sequel. It is consistent with the structural formulas and numerical
behaviour of \cite{Xiao2001,OsipovRokhlinXiao}, but is not claimed here as a
known theorem.
\end{conjecture}

\begin{remark}[Why this is an implication framework]
Conjectures~\ref{conj:pswf_product_tail} and~\ref{conj:xry_stability}
isolate the two unresolved analytical inputs used in this paper.
Accordingly, the present work should be read as an implication framework
for the DSTP of \cite{PaperI}, not as an
unconditional verification of that compatibility condition.
The scope of the implication is now restricted to the bulk and off-diagonal
regimes; the edge is excluded by Conjecture~\ref{conj:pswf_product_tail}(iii).
\end{remark}

%% ==================================================================
\section{Conditional Verification of the DSTP}\label{sec:verification}
%% ==================================================================

\begin{proposition}[Structured defect bound]\label{prop:structured_defect}
Assume Conjecture~\ref{conj:xry_stability}. Then for all $0\le m,n\le N-1$,
\[
  |E_{mn}|
  \le \CXRY\,\|(I-\PN)f_{mn}\|_{L^2([-T,T])}
      + C\,e^{-\alpha N}\,\norm{f_{mn}}_{L^2([-T,T])}.
\]
\end{proposition}
\begin{proof}
The exact identity
$E_{mn}=R_{mn}^{\mathrm{quad}}=Q_N(f_{mn})-\int_{-T}^T f_{mn}(x)\,dx$
from \cite{PaperI} and Conjecture~\ref{conj:xry_stability}.
\end{proof}

\begin{proposition}[Conditional exponential quadrature defect bound
--- bulk/off-diagonal]\label{prop:defect_bound}
Assume part~(i) or part~(ii) of
Conjecture~\ref{conj:pswf_product_tail},
Conjecture~\ref{conj:xry_stability}, and Assumption~\ref{ass:pswf_sup},
for indices in the respective regime.
Then
\[
  |E_{mn}|\le C\,e^{-\alpha N}
  \quad\text{(bulk, part~(i))},
  \qquad
  |E_{mn}|\le C_p(1+|m-n|)^{-p}
  \quad\text{(off-diagonal, part~(ii))}.
\]
\end{proposition}
\begin{proof}
Insert the respective tail bound into
Proposition~\ref{prop:structured_defect} and use
Lemma~\ref{lem:product_norm} for the $L^2$ factor.
\end{proof}

\begin{remark}[No gradual interpolation in the present framework]
Within the present argument, there is no gradual weakening of the
conclusion obtainable from partial versions of the two conjectures. The
exponential defect bound follows only once both conjectural inputs are
available in the stated form within the relevant regime.
\end{remark}

\begin{proposition}[Partial compatibility under bulk/off-diagonal control]
\label{prop:partial_compatibility}
Assume Conjecture~\ref{conj:xry_stability}, Assumption~\ref{ass:pswf_sup},
and the following partial PSWF tail information:
\begin{enumerate}[(a)]
\item \emph{Unconditional $L^2$ bound} (Theorem~1 of \cite{PaperIII}):
  for all $0\le m,n\le N$,
  \[
    \|(I-\PN)f_{mn}\|_{L^2([-T,T])} \le C T^{1/2}.
  \]
  This bound is unconditional but does not decay in $|m-n|$.
\item \emph{Conditional algebraic bound} (Theorem~2 of \cite{PaperIII},
  under Assumption~3.1 therein): for $|m-n|\ge\delta N$,
  \[
    \|(I-\PN)f_{mn}\|_{L^2([-T,T])}
    \le C_p(1+|m-n|)^{-p}.
  \]
\end{enumerate}
From~(a), the unconditional defect estimate is
$|E_{mn}| \le C'T^{1/2} + Ce^{-\alpha N}$,
which does not decay and gives no banded structure.
From~(b), the conditional defect estimate is
\[
  |E_{mn}|\le C_p'(1+|m-n|)^{-p}+C e^{-\alpha N}
\]
in the off-diagonal regime, giving a banded-type compatibility statement
consistent with the DSTP of \cite{PaperI}, away from the edge.
\end{proposition}

\begin{remark}[Unconditional bound is insufficient for frame stability]
\label{rem:unconditional-insufficient}
The unconditional $T^{1/2}$ bound from part~(a) of
Proposition~\ref{prop:partial_compatibility} does not exhibit decay
in $|m-n|$ and therefore does not imply the banded structure required
by the DSTP of \cite{PaperI}.
Frame stability in the off-diagonal regime requires the conditional
algebraic decay of part~(b), i.e., the spectral localization input
of Assumption~3.1 of \cite{PaperIII}.
\end{remark}

\begin{lemma}[Exponential domination of the bulk error scale]
\label{lem:bulk_scale}
Assume $N=\gamma c$ with fixed $\gamma>0$. Then for every $\alpha>0$ there
exists $C_{\alpha,\gamma}>0$ such that
$e^{-\alpha N}\le C_{\alpha,\gamma} c^{-1/3}$
for all sufficiently large $c$.
\end{lemma}
\begin{proof}
$e^{-\alpha N}=e^{-\alpha\gamma c}$ decays faster than any power of $c$.
\end{proof}

\begin{remark}[Uniform domination over the index range]\label{rem:index_domination}
For every $p\ge1$ and sufficiently large $N$, $e^{-\alpha N}\le C_p(1+N)^{-p}$.
This is the precise sense in which an exponential bound in $N$ dominates
the algebraic defect profile required by the DSTP of \cite{PaperI}.
\end{remark}

\begin{theorem}[Conditional verification theorem --- bulk
regime]\label{thm:compatibility_verified}
Assume part~(i) of Conjecture~\ref{conj:pswf_product_tail},
Conjecture~\ref{conj:xry_stability}, and Assumptions~\ref{ass:pswf_sup},
\ref{ass:node_sep}, \ref{ass:weights}. Let $(\mathbf{p},\mathbf{w})$
be the XRY PSWF-Gauss rule of order $N=\gamma c$ with fixed $\gamma>0$.
For indices $m,n\le\gamma N$, every $p\ge1$, there exist $C_p,C>0$ such that
\[
  |E_{mn}|\le C_p(1+|m-n|)^{-p}+Cc^{-1/3}.
\]
Hence the DSTP of \cite{PaperI} holds conditionally for
PSWF-Gauss sampling in the bulk regime.
The edge regime $m,n\sim N$ is excluded;
see Conjecture~\ref{conj:pswf_product_tail}(iii).
\end{theorem}
\begin{proof}
By Proposition~\ref{prop:defect_bound} (part~(i)),
$|E_{mn}|\le Ce^{-\alpha N}$.
Lemma~\ref{lem:bulk_scale} and Remark~\ref{rem:index_domination}
absorb this into the stated form.
\end{proof}

\begin{remark}[What has been achieved]
The preceding theorem does not close the construction problem of
\cite{PaperI} unconditionally, nor globally: the edge regime
$m,n\sim N$ is excluded by the proved obstruction of
Conjecture~\ref{conj:pswf_product_tail}(iii).
What it provides is a precise conditional implication in the bulk:
once Conjectures~\ref{conj:pswf_product_tail}(i) and
\ref{conj:xry_stability} are established, the DSTP
of \cite{PaperI} follows for XRY PSWF-Gauss quadrature
in the bulk regime.
\end{remark}

\begin{corollary}[Conditional defect norm bound]\label{cor:defect_norm}
Under the assumptions of Theorem~\ref{thm:compatibility_verified},
restricting to $m,n\le\gamma N$:
$\norm{\EN}_2\le\norm{\EN}_F\le N C_* e^{-\alpha N}\to0$.
\end{corollary}
\begin{proof}
$N^2$ entries each bounded by $C_*e^{-\alpha N}$;
$\norm{\EN}_F\le NC_*e^{-\alpha N}$.
\end{proof}

%% ==================================================================
\section{Conditional Frame Stability for PSWF-Gauss Sampling}\label{sec:stability}
%% ==================================================================

\begin{theorem}[Conditional frame stability theorem]\label{thm:frame_stability}
Under the assumptions of Theorem~\ref{thm:compatibility_verified}
and $\gamma<2/\pi$, once $NC_*e^{-\alpha N}<\lambda_{N-1}(c)$,
$\{(p_j,w_j)\}_{j=1}^N$ is a weighted frame for $V_{N,c}$ with bounds
\[
  A(N,c)\,\norm{f}_{L^2}^2
  \le \sum_{j=1}^N w_j|f(p_j)|^2
  \le B(N,c)\,\norm{f}_{L^2}^2,
  \quad f\in V_{N,c},
\]
where $A(N,c)=\lambda_{N-1}(c)-NC_*e^{-\alpha N}>0$,
$B(N,c)=\lambda_0(c)+NC_*e^{-\alpha N}$.
\end{theorem}
\begin{proof}
Weyl perturbation: $\lambda_{\min}(\GVM)\ge\lambda_{N-1}(c)-\norm{\EN}_2
\ge A(N,c)>0$ by Corollary~\ref{cor:defect_norm}.
\end{proof}

\begin{remark}[Critical density regime]\label{rem:critical}
Near the critical density $N\to(2/\pi)c$,
$\lambda_{N-1}(c)\to0$ and the stability condition may fail;
see Open Problem~\ref{op:critical}.
\end{remark}

%% ==================================================================
\section{Open Problems}\label{sec:open}
%% ==================================================================

\begin{problem}[Derive the bulk PSWF product spectral tail conjecture]
\label{prob:bulkconj}
Prove Conjecture~\ref{conj:pswf_product_tail}(i) rigorously.
Natural route: exponential WKB decay of PSWF Fourier tails near $|\xi|=\omega$
for $n\le\gamma N$; see Problem~1 of \cite{PaperIII}.
\end{problem}

\begin{problem}[Derive the off-diagonal spectral localization]
Prove Conjecture~\ref{conj:pswf_product_tail}(ii) rigorously,
i.e., Assumption~3.1 of \cite{PaperIII}.
The energy equidistribution question of Problem~4.2 of \cite{PaperIII}
is the identified analytic bottleneck.
\end{problem}

\begin{problem}[Derive the XRY stability conjecture]
Prove Conjecture~\ref{conj:xry_stability} rigorously from \cite{Xiao2001,OsipovRokhlinXiao}.
\end{problem}

\begin{problem}[Node separation and weight bounds]
Prove Assumptions~\ref{ass:node_sep} and~\ref{ass:weights} quantitatively.
\end{problem}

\begin{problem}[Critical density regime]\label{op:critical}
Analyse frame stability near $N\to(2/\pi)c$.
\end{problem}

\begin{problem}[Edge regime via doubled prolate subspace]
Since the edge obstruction is proved, investigate whether the
doubled projector $P_{2N,2c}$ can substitute for $\PN$ in the
DSTP argument; see Problem~5 of \cite{PaperIII}.
\end{problem}

\begin{problem}[Alternative compatible quadrature classes]
Characterise quadrature rules beyond the XRY PSWF-Gauss class satisfying
the DSTP of \cite{PaperI}.
\end{problem}

%% ==================================================================
\begin{thebibliography}{99}

\bibitem{PaperI}
Anonymous Author,
\emph{Frame Coercivity of Discrete Prolate Sampling Operators
under Quadrature Compatibility: Defect Decomposition and a
Bulk--Turning Decay Mechanism},
preprint, 2026.

\bibitem{PaperIII}
Anonymous Author,
\emph{PSWF Product Spectral Tail Estimates: Bulk and Off-Diagonal Regimes,
and an Obstruction at the Spectral Edge
(Paper~III in the Prolate--Weil Program)},
preprint, 2026.

\bibitem{OsipovRokhlinXiao}
A.~Osipov, V.~Rokhlin, and H.~Xiao,
\textit{Prolate Spheroidal Wave Functions of Order Zero}, Springer, 2013.

\bibitem{Slepian1978}
D.~Slepian,
\textit{Prolate spheroidal wave functions, Fourier analysis, and
uncertainty~-- V},
Bell Syst.\ Tech.\ J.\ \textbf{57} (1978), 1371--1430.

\bibitem{Xiao2001}
H.~Xiao, V.~Rokhlin, and N.~Yarvin,
\textit{Prolate spheroidal wave functions, quadrature, and interpolation},
Inverse Problems \textbf{17} (2001), 805--838.

\end{thebibliography}

\end{document}
